\title{คาบเรียนที่ 7: การประยุกต์ของอนุพันธ์ (1/2)}
\frame{\titlepage}

\begin{frame}
ในที่นี้ เราจะนำค่าอนุพันธที่ได้มาประยุกต์ใช้ โดยจะให้ความหมายกับค่าของ $\dyx$ ที่หามาได้ ในหัวข้อต่าง ๆ ดังต่อไปนี้
\tableofcontents
\end{frame}


%%%%%%%%%%%%%%
\section[อันดับสูง]{อนุพันธ์อันดับสูง}
\frame{\sectionpage}
%%%%%%%%%%%%%%

\begin{frame}
\frametitle{อนุพันธ์อันดับสอง}
เราสามารถหาอนุพันธ์ต่อเนื่องได้ ซึ่งจะมีความหมายแตกต่างกันไปแล้วแต่หัวข้อที่ใช้ศึกษา ในขั้นนี้เราเพียงให้คำนิยามและแสดงตัวอย่างเท่านั้น
\begin{defn}{}{}
กำหนด $y = f(x)$ ถ้า $f\,'(x)$ เป็นฟังก์ชันที่หาอนุพันธ์ได้ \pa เราจะเรียก \textbf{อนุพันธ์อันดับสอง} ของ $f$ เทียบกับ $x$ ว่า

\[ \dx[2] f(x) = \dx \left( \dx f(x) \right) = \dx \left( f\,'(x) \right)\]

\pa โดยเราสามารถเขียนในรูปแบบอื่นได้ดังนี้
\[\dx[2] f(x) = f\,''(x) = f^{\,(2)}(x)\]
\end{defn}
\end{frame}


\begin{frame}[t]
\begin{ex}
กำหนดให้ $f(x) = x^3$ จงหาอนุพันธ์อันดับสอง $f\,''(x)$ \pa
\end{ex}
\begin{sol} 
\begin{eqnarray*}
f\,''(x) = \dx[2] f(x) \pa &=& \dx \left( \dx x^3 \right)  \pa\\
&=&  \dx (3x^2) \pa \\
&=& 3 \dx x^2 \pa \\
&=& 3(2x) \pa = 6x
\end{eqnarray*}
\end{sol}
\end{frame}


\begin{frame}[t]
\begin{ex}
กำหนดให้ $f(x) = \frac{1}{x^2}$ จงหา $f\,''(x)$ \pa
\end{ex}
\begin{sol}
\begin{eqnarray*}
f\,''(x) = \dx[2] f(x) \pa &=& \dx \left( \dx x^{-2} \right) \pa \\
&=&  \dx \left( -2(x^{-3}) \right) \pa \\
&=& (-2) \dx x^{-3} \pa \\
&=& (-2)(-3x^{-4}) \pa = 6x^{-4} 
\end{eqnarray*}
\end{sol}
\end{frame}


\begin{frame}[t]
\begin{ex}
กำหนดให้ $f(x) = xe^x$ จงหา $f\,''(x)$ \pa
\end{ex}
\begin{sol}
\begin{eqnarray*}
f\'(x) &=& \dx xe^x \\
&=& x \dx e^x + e^x \dx x \\
&=& xe^x+e^x \\
f\,''(x) &=& \dx \left( xe^x + e^x \right) \pa \\
&=&  \left(\dx xe^x\right) + \left(\dx  e^x\right) \\
&=& (xe^x + e^x) + e^x \\
&=& xe^x + 2e^x
\end{eqnarray*}
\end{sol}
\end{frame}

\begin{frame}[t]
\begin{ex}
กำหนดให้ $f(x) = \sin x + \cos x$ จงหา $f\,''(x)$ \pa
\end{ex}
\begin{sol}
\begin{eqnarray*}
f\,'(x) &=&  \dx (\sin x + \cos x) \\
&=& \left(\dx \sin x\right) + \left(\dx \cos x \right) \mathnote{สมบัติอนุพันธ์ของผลบวกฟังก์ชัน} \\
&=& \cos x - \sin x \\
f\,''(x) = \dx f\,'(x) \pa &=&  \dx \left( \cos x - \sin x \right) \pa \\
&=& \left(\dx \cos x\right) - \left(\dx \sin x\right) \\
\pa &=& -\sin x -\cos x
\end{eqnarray*}
\end{sol}
\end{frame}

\begin{frame}
\frametitle{อนุพันธ์ของฟังก์ชันอันดับสูง}
\begin{defn}{}{}
\textbf{อนุพันธ์อันดับ} $n$ ของ $f$ เทียบ $x$ \pa นิยามโดย
\begin{eq}
\dx[n] f(x) \pa &=& \dx \left( f^{\,(n-1)(x)} \right) \\
\pa &=& \dx \left( \dx \left( \cdots \left( \dx f(x) \right) \cdots \right) \right) \nonumber
\end{eq}
\pa โดยมีพจน์ $\dx$ ทั้งหมด $n$ พจน์ \pa ซึ่งหมายถึงการหาอนุพันธ์ต่อเนื่อง รวมทั้งหมด $n$ ครั้ง ซึ่งเราสามารถเขียนแบบย่อได้เป็น 
\[ \dx[n] f(x) = f^{\,(n)}(x)\]
\end{defn}
\end{frame}


\begin{frame}[t]
\begin{ex}
กำหนดให้ $f(x) = x^4$ จงหาอนุพันธ์อันดับ 4 (นั่นคือ $f\,^{(4)}(x)$)
\end{ex}
\begin{sol} เราหาอนุพันธ์อันดับ 1,2,3 และ 4 ไล่ขึ้นไปทีละอันดับได้ดังนี้
\begin{eqnarray*}
\red{f\,'(x)} &=& \dx (x^4) \pa = 4x^3 \\
\green{f\,''(x)} = \dx \red{f\,'(x)} \pa &=& \dx (4x^3) \pa = 4(3x^2) \pa = 12x^2 \\
\blue{f\,'''(x)} = \dx \green{f\,''(x)} \pa &=& \dx (12x^2) \pa = 12(2x) \pa = 24x \\
f\,^{(4)} = \dx \blue{f\,'''(x)} \pa &=& \dx (24x) \pa = 24
\end{eqnarray*}
ฉะนั้น $f\,^{(4)}(x) = 24$
\end{sol}
\end{frame}


\begin{frame}[t]
\begin{ex}
กำหนดให้ $f(x) = x^5+x^2$ จงหา $f\,'''(x)$
\end{ex}
\begin{sol} เราหาอนุพันธ์อันดับ 1,2 และ 3 ไล่ขึ้นไปทีละอันดับได้ดังนี้
\begin{eqnarray*}
\red{f\,'(x)} &=& \dx (x^5 + x^2) \pa = \left( \dx x^5 \right) + \left( \dx x^2 \right)  \\
\pa &=& 5x^4 + 2x \\
\pa \blue{f\,''(x)} = \dx \red{f\,'(x)} \pa &=& \dx (5x^4 + 2x) \pa = \left( \dx 5x^4 \right) + \left( \dx 2x \right) \\
\pa &=& 5(4x^3) + 2(1) \pa = 20x^3 + 2 \\
f\,'''(x) = \dx \blue{f\,''(x)} \pa &=& \dx (20x^3 + 2) = \left( \dx 20x^3 \right) + \left( \dx 2 \right)\\
\pa &=& 20(3x^2) + 0 \pa = 60x^2
\end{eqnarray*}
ฉะนั้น $f\,'''(x) = 60x^2$
\end{sol}
\end{frame}


\begin{frame}[t]
\begin{ex}
กำหนดให้ $f(x) = e^{2x}$ จงหาอนุพันธ์อันดับ 3 (นั่นคือ $f\,^{(3)}(x)$)
\end{ex}
\begin{sol} เราหาอนุพันธ์อันดับ 1,2 และ 3 ไล่ขึ้นไปทีละอันดับได้ดังนี้
\begin{eqnarray*}
\red{f\,'(x)} &=& \dx (e^{2x})= e^{2x} \dx (2x) \mathnote{จากสูตร $\dx e^u = e^u \dx u$}\\
&=& 2e^{2x} \\
\green{f\,''(x)} &=& \dx \red{f\,'(x)} = \dx (2e^{2x}) = 2e^{2x} \dx (2x) = 4e^{2x} \\
f\,'''(x) &=& \dx \green{f\,''(x)} = \dx (4e^{2x}) = 4e^{2x} \dx 2x = 8e^{2x}
\end{eqnarray*}
ฉะนั้น $f\,^{(3)}(x) = 8e^{2x}$
\end{sol}
\end{frame}

\begin{frame}
\frametitle{หมายเหตุ}
\begin{remark}{}
การหาอนุพันธ์อันดับสูง แนะนำให้ใช้สัญลักษณ์ $\dx[n] f(x)$ แทน $f\,^{(n)}$ แต่ถ้าต้องการใช้ตัวย่อ ให้ระวังการใช้ดังนี้
\begin{itemize}
\item เราสามารถใชัตัวย่อ $f\,', f\,''$ และ $f\,'''$ แทน $f\,^{(1)}, f\,^{(2)}$ และ $f\,^{(3)}$ ได้ 
\item สำหรับอนุพันธ์อันดับที่ 4 เป็นต้นไป เราจะใช้สัญลักษณ์ $f\,^{(n)}(x)$ ตามนิยามเสมอ ไม่ใช้ $'$ เนื่องจากสามารถนับได้ยากว่ามีสัญลักษณ์นั้นทั้งหมดกี่ตัว (เช่น $f\,'''''(x)$ มีสัญลักษณ์ $'$ อยู่กี่ตัว?)
\item เราจำเป็นจะต้องใส่วงเล็บรอบตัวเลข $(n)$ เพราะถ้าไม่ใส่แล้วจะกลายเป็น $f\,^{n}(x)$ ซึ่งอาจจะหมายความว่านำฟังก์ชันไปยกกำลัง $f\,^n(x) = (f (x))^n$ ก็ได้
\end{itemize}
\end{remark}
\end{frame}

%%%%%%%%%%%%%%
\section{อนุพันธ์ในทางเรขาคณิต}
\frame{\sectionpage}
%%%%%%%%%%%%%%

\begin{frame}
เราได้ศึกษานิยามของอนุพันธ์ ซึ่งเท่ากับ
\[\dyx = \limh \frac{f(x+h)-f(x)}{h}\]
โดยเราให้ความหมายว่าเป็นอัตราความเปลี่ยนแปลงของตัวแปรตาม $y$ เทียบกับตัวแปรอิสระ $x$
\end{frame}

\begin{frame}
\frametitle{ทบทวนเส้นตรง}
เราได้เคยศึกษาสมการเส้นตรงมาแล้ว ดังนี้
\begin{thm}{}{}
สมการเส้นตรงที่ผ่านจุด $(x_0, y_0)$ \pa และมีความชันเท่ากับ $m$ \pa คือ
\begin{equation} 
y-y_0 = m(x-x_0) \label{eq:applied_basic_line}
\end{equation} 
\end{thm}
\addfig{0.3}{diff_applied_line}{สมการเส้นตรงที่ผ่านจุด $(x_0, y_0)$ และมีค่าชันเท่ากับ $m$}{}
\end{frame}

\begin{frame}
\frametitle{สมการเส้นสัมผัสเส้นโค้ง}
กำหนดกราฟ $y=f(x)$ และเราต้องการหาสมการเส้นสัมผัสเส้นโค้งเมื่อกำหนด $x=x_0$ เราจะหาค่า $y_0$ และ $m$ ได้ดังนี้
\begin{itemize}
\item จากกราฟที่กำหนดให้ เราทราบว่า $y_0 = f(x_0)$ นั่นคือจุด $(x_0, f(x_0))$ เป็นจุดที่เส้นสัมผัสผ่าน
\item จากนิยามของอนุพันธ์ เราทราบว่าความชันของเส้นสัมผัสเส้นโค้งเมื่อ $x=x_0$ คือ $m = f\,'(x_0)$
\end{itemize}
\addfig{0.3}{diff_applied_tangent_line}{เส้นสัมผัสเส้นโค้ง $y=f(x)$ เมื่อ $x = x_0$}{}
\end{frame}

\begin{frame}
ฉะนั้น เมื่อนำจุด $(x_0, f(x_0))$ และความชัน $m = f\,'(x_0)$ ไปใส่ในสมการ \eqref{eq:applied_basic_line} จะสามารถสรุปเป็นทฤษฎีบทได้ดังนี้
\begin{thm}{}{}
กำหนดฟังก์ชัน $f(x)$ ที่หาอนุพันธ์ได้ที่จุด $x=x_0$ \pa แล้วสมการเส้นสัมผัสเส้นโค้ง $y = f(x)$ ที่จุด $(x_0, f(x_0))$ \pa คือ
\begin{equation} \label{eq:applied_line}
 y - f(x_0) = f\,'(x_0)(x-x_0)
 \end{equation}
\end{thm}
\end{frame}


\begin{frame}[t]
\begin{ex}
จงหาสมการเส้นสัมผัสเส้นโค้ง $y = x^3$ เมื่อ $x=2$
\end{ex}
\begin{sol} กำหนดให้ $f(x) = x^3$
\begin{itemize}
\item เนื่องจาก $f(\red{2}) = \red{2}^3 = \blue{8}$ \pa ดังนั้นจุด $(\red{2}, f(\red{2})) = (\red{2},\blue{8})$ อยู่บนเส้นโค้ง $y = f(x)$ \pa
\item ความชันของเส้นสัมผัสเส้นโค้งที่จุด $(2, 8)$ คืออนุพันธ์ของฟังก์ชัน \pa ซึ่งเราหาอนุพันธ์ได้เป็น $f\,'(x) = 3x^2$ \pa ฉะนัั้นเมื่อแทนค่า $\red{x=2}$ ลงไปจะได้ $f\,'(\red{2}) = 3(\red{2})^2 = \green{12}$ \pa
\end{itemize}
ฉะนั้น จากสมการ \eqref{eq:applied_line} สมการเส้นสัมผัสเส้นโค้ง $y = f(x)$ ที่จุด $(1,2)$ คือ \pa
\begin{eqnarray*}
y - f(x_0) &=& f\,'(x_0) (x-x_0) \\
y - \blue{8} \pa &=& \green{12}(x - \red{2}) \pa \\
y &=& 12x-16
\end{eqnarray*}
\end{sol}
\end{frame}

\begin{frame}
\addfig{0.8}{diff_applied_tangent_ex1}{ตัวอย่างที่ \theex: สมการเส้นสัมผัส $y=12x-16$}{}
\end{frame}


\begin{frame}[t]
\begin{ex}
จงหาสมการเส้นสัมผัสเส้นโค้ง $y = 5x^2 - 3$ เมื่อ $x=1$
\end{ex}
\begin{sol} กำหนดให้ $f(x) = 5x^2-3$
\begin{itemize}
\item เนื่องจาก $f(1) = 5(1^2) - 3 = 2$ \pa ดังนั้นจุด $(1, f(1)) = (1,2)$ อยู่บนเส้นโค้ง $y = f(x)$ \pa
\item ความชันของเส้นสัมผัสเส้นโค้งที่จุด $(1, 2)$ คืออนุพันธ์ของฟังก์ชัน \pa ซึ่งเรารู้ว่า $f\,'(x) = 10x$ \pa ฉะนัั้น $f\,'(1) = 10(1) = 10$ \pa
\end{itemize}
ฉะนั้น จากสมการ \eqref{eq:applied_line} สมการเส้นสัมผัสเส้นโค้ง $y = f(x)$ ที่จุด $(1,2)$ คือ \pa
\begin{eqnarray*}
y - f(x_0) &=& f\,'(x_0) (x-x_0) \pa \\
y - 2 \pa &=& 10(x-1) \pa \\
y &=& 10x-8
\end{eqnarray*}
\end{sol}
\end{frame}

\begin{frame}
\addfig{0.8}{diff_applied_tangent_ex2}{ตัวอย่างที่ \theex: สมการเส้นสัมผัส $y=12x-16$}{}
\end{frame}


\begin{frame}[t]
\begin{ex}
จงหาสมการเส้นสัมผัสเส้นโค้ง $y = \ln (x+3)$ เมื่อ $x=5$
\end{ex}
\begin{sol} กำหนดให้ $f(x) = \ln (x+3)$
\begin{itemize}
\item เนื่องจาก $f(5) = \ln (5+3) = \ln 8$ \pa ดังนั้นจุด $(5, f(5)) = (5, \ln 8)$ อยู่บนเส้นโค้ง $y = f(x)$ \pa
\item ความชันของเส้นสัมผัสเส้นโค้งที่จุด $(5,\ln 8)$ คืออนุพันธ์ของฟังก์ชัน \pa ซึ่งเราหาอนุพันธ์ได้เป็น $f\,'(x) = \tfrac{1}{x+3}$ \pa ฉะนัั้นเมื่อแทนค่า $x=5$ ลงไปจะได้ $f\,'(5) = \tfrac{1}{5+3} = \tfrac{1}{8}$ \pa
\end{itemize}
ฉะนั้น จากสมการ \eqref{eq:applied_line} สมการเส้นสัมผัสเส้นโค้ง $y = f(x)$ ที่จุด $(5, \ln 8)$ คือ \pa
\begin{eqnarray*}
y - f(x_0) &=& f\,'(x_0) (x-x_0) \\
y - \ln 8 \pa &=& \tfrac{1}{8}(x - 5) \pa \\
y &=& \tfrac{x}{8} - \tfrac{5}{8} + \ln 8
\end{eqnarray*}
\end{sol}
\end{frame}

\begin{frame}
\addfig{0.8}{diff_applied_tangent_ex3}{ตัวอย่างที่ \theex: สมการเส้นสัมผัส $y=10x-8$}{}
\end{frame}


\begin{frame}[t]
\begin{ex}
ถ้าเส้นโค้ง $y = x^2+4$ สัมผัสกับเส้นตรง $y = mx$ เมื่อ $m$ เป็นจำนวนจริงที่ยังไม่ทราบค่า แล้วจงหาจุดสัมผัสและค่า $m$
\end{ex}
\begin{sol}
กำหนด $f_1(x) = x^2+4$ และ $f_2(x) = mx$ และให้เส้นตรงสัมผัสกับเส้นโค้งที่จุด $(x_0, y_0)$ แล้วเราทราบว่า
\begin{enumerate}
\item จุดที่สัมผัสต้องอยู่บนเส้นโค้งและเส้นตรง นั่นคือ
\[ f_1(x_0) = f_2(x_0) \]
\item ความชันของเส้นโค้ง ณ จุดนั้น ต้องเท่ากับความชันของเส้นตรง ณ จุดนั้น นั่นคือ
\[ f_1\,'(x_0) = f_2\,'(x_0) \]
\end{enumerate}
เราสามารถใช้สองสมการนี้ในการหาค่า $m$ และ $(x_0, y_0)$ ได้ดังนี้
\end{sol}
\end{frame}

\begin{frame} 
\begin{sol} (ต่อ)
จากสมการที่สอง จะได้
\begin{eq*}
f_1\,'(x_0) &=& f_2\,'(x_0) \\
\pa 2x_0 &=& m
\end{eq*}
เมื่อได้ค่า $m$ แล้ว ให้นำจุดนี้ไปแทนในสมการแรก จะได้
\begin{eq*}
f_1(x_0) &=& f_2(x_0) \\
x_0^2 +4 &=& mx_0 \\
&=& (2x_0)x_0 \\
x_0^2 &=& 4 \\
x_0 &=& -2, 2
\end{eq*}
ฉะนั้น เราได้สองคำตอบดังนี้
\begin{itemize}
\item จุดสัมผัสคือจุด $(-2, f_1(-2)) = (2, 8)$ โดยมีค่าชันเท่ากับ $m = 2(-2) = -4$
\item จุดสัมผัสคือจุด $(2, f_1(2)) = (2, 8)$ โดยมีค่าชันเท่ากับ $m = 2(2) = 4$
\end{itemize}
\end{sol}
\end{frame}

\begin{frame}
\addfig{0.8}{diff_applied_tangent_ex4}{ตัวอย่างที่ \theex: สมการเส้นสัมผัส $y=4x$ และ $y=-4x$}{}
\end{frame}

%%%%%%%%%%%%%%
\section{ฟังก์ชันลด ฟังก์ชันเพิ่ม และความเว้า}
\frame{\sectionpage}
%%%%%%%%%%%%%%

\begin{frame}
\addfig{0.8}{diff_applied_extrema_graph}{ตัวอย่างกราฟราคาหุ้น \href{https://www.pexels.com/photo/blue-and-yellow-graph-on-stock-market-monitor-159888/}{(ที่มา)}}{}
\end{frame}

\begin{frame}
\frametitle{ฟังก์ชันเพิ่ม ฟังก์ชันลด และฟังก์ชันคงค่า}
\begin{defn}{}{}
กำหนดให้ $f$ เป็นฟังก์ชันนิยามบนช่วงปิด $[a,b]$ 
\begin{enumerate}
\item ถ้า $f(x_1) < f(x_2)$ สำหรับทุก $x_1 < x_2$ ใน $[a,b]$ \pa จะเรียก $f$ ว่า \textbf{ฟังก์ชันเพิ่ม} \pa
\item ถ้า $f(x_1) > f(x_2)$ สำหรับทุก $x_1 < x_2$ ใน $[a,b]$ \pa จะเรียก $f$ ว่า \textbf{ฟังก์ชันลด} \pa
\item ถ้า $f(x_1) = f(x_2)$ สำหรับทุก $x_1 < x_2$ ใน $[a,b]$ \pa จะเรียก $f$ ว่า \textbf{ฟังก์ชันคงค่า} \pa
\end{enumerate}
\end{defn}
\end{frame}

\begin{frame}
\addfig{1}{diff_applied_increasing}{(1) กราฟของฟังก์ชันเพิ่ม (2) กราฟของฟังก์ชันลด (3) กราฟของฟังก์ชันคงค่า}{}
\end{frame}


\begin{frame}[t]
\begin{ex}
\begin{columns}
\begin{column}{0.6\textwidth}
จงพิจารณาว่าฟังก์ชัน $f(x) = x^3-3x$ เป็นฟังก์ชันเพิ่มในช่วงใด และเป็นฟังก์ชันลดในช่วงใด โดยการพิจารณาจากกราฟ
\end{column}
\begin{column}{0.4\textwidth}
\addfig{1}{diff_applied_increasing_ex1}{กราฟ $y = x^3-3x$}{}
\end{column}
\end{columns}
\end{ex}
\begin{sol}
จากกราฟจะได้ว่า $f(x)$ เป็นฟังก์ชันเพิ่มในช่วง $(-\infty, -1]$ และ $[1, \infty)$ และเป็นฟังก์ชันลดในช่วง $[-1,1]$

ถ้ากำหนดช่วง $[-2,2]$ มาให้พิจารณา เราจะสรุปได้ว่า $f(x)$ ไม่เป็นทั้งฟังก์ชันเพิ่มและฟังก์ชันลดในช่วง $[-2,2]$
\end{sol}
\end{frame}

\begin{frame}[t]
\frametitle{ความสัมพันธ์ระหว่างอนุพันธ์และพฤติกรรมของฟังก์ชัน}
พิจารณาจุด $(x_1,f(x_1))$ และจุด $(x_2, f(x_2))$ เมื่อ $x_1 < x_2$
\begin{slide}
\begin{itemize}
\item ถ้า $f(x)$ เป็นฟังก์ชันเพิ่มบนช่วง $[a,b]$ จะได้ว่า $f(x_2) - f(x_1) > 0$ และ $x_2 - x_1 > 0$ ฉะนั้น 
\[\frac{f(x_2)-f(x_1)}{x_2-x_1} > 0\]
\end{itemize}
ซึ่งจะได้ว่า $f\,'(x) \geq 0$ สำหรับทุกค่า $x$ ในช่วง $[a,b]$
\end{slide}
\addfig{0.3}{diff_applied_derivative_increasing}{อนุพันธ์ในช่วงของฟังก์ชันเพิ่ม}{}
\end{frame}

\begin{frame}
\begin{slide}
\begin{itemize}
\item ในทำนองเดียวกัน ถ้า $f(x)$ เป็นฟังก์ชันลดบนช่วง $[a,b]$ จะได้ว่า $f(x_2) - f(x_1) < 0$ และ $x_2 - x_1 > 0$ ฉะนั้น 
\[\frac{f(x_2)-f(x_1)}{x_2-x_2} < 0\]
\end{itemize}
ซึ่งจะได้ว่า $f\,'(x) \leq 0$ สำหรับทุกค่า $x$ ในช่วง $[a,b]$
\end{slide}
\addfig{0.3}{diff_applied_derivative_decreasing}{อนุพันธ์ในช่วงของฟังก์ชันลด}{}
\end{frame}

\begin{frame}
\begin{slide}
\begin{itemize}
\item สุดท้าย ถ้า $f(x)$ เป็นฟังก์ชันคงค่าบนช่วง $[a,b]$ จะได้ว่า $f(x_2) - f(x_1) = 0$ ฉะนั้น 
\[\frac{f(x_2) - f(x_1)}{x_2 - x_1} = 0\]
\end{itemize}
ซึ่งจะได้ว่า $f\,'(x) = 0$ สำหรับทุกค่า $x$ ในช่วง $[a,b]$
\end{slide}
\addfig{0.5}{diff_applied_derivative_constant}{อนุพันธ์ในช่วงของฟังก์ชันคงค่า}{}
\end{frame}


\begin{frame}[t]
\begin{ex}
จงพิจารณาว่าฟังก์ชัน $f(x) = x^3-3x$ เป็นฟังก์ชันเพิ่มในช่วงใด และเป็นฟังก์ชันลดในช่วงใด โดยใช้อนุพันธ์
\end{ex}
\begin{sol}
หาอนุพันธ์ได้  $f\,'(x) = 3x^2-3$ จากนั้นพิจารณาดังนี้
\begin{itemize}
\item $f\,'(x) > 0$ เมื่อ 
\begin{eq*}
3x^2 -3 &>& 0 \\
(x+1)(x-1) & >& 0 \\
x > 1 &\text{ หรือ }& x < -1
\end{eq*}
นั่นคือ $f(x)$ เป็นฟังก์ชันเพิ่มในช่วง $(-\infty, -1)$ และ $(1,\infty)$
\item $f\,'(x) < 0$ เมื่อ $-1 < x < 1$ ฉะนั้น $f(x)$ เป็นฟังก์ชันลดในช่วง $(-1,1)$
\item $f\,'(x) = 0$ เมือ $x=-1$ หรือ $x=1$ ฉะนั้นไม่มีช่วงที่ $f(x)$ เป็นฟังก์ชันค่าคงตัว
\end{itemize}
\end{sol}
\end{frame}

\begin{frame}
\begin{thm}{}{}
กำหนดให้ $f$ เป็นฟังก์ชันนิยามบนช่วงปิด $[a,b]$ 
\begin{enumerate}
\item ถ้า $f\,'(x) > 0$ สำหรับทุก $x_1 < x_2$ ใน $(a,b)$ แล้ว $f$ เป็นฟังก์ชันเพิ่มบน $[a,b]$ \pa
\item ถ้า $f\,'(x) < 0$ สำหรับทุก $x_1 < x_2$ ใน $(a,b)$ แล้ว $f$ เป็นฟังก์ชันลดบน $[a,b]$  \pa
\item ถ้า $f\,'(x) = 0$ สำหรับทุก $x_1 < x_2$ ใน $(a,b)$ แล้ว $f$ เป็นฟังก์ชันคงค่าบน $[a,b]$  \pa
\end{enumerate}
\end{thm}
\end{frame}

\begin{frame}
\frametitle{ความหมายของอนุพันธ์อันดับสอง}
พิจารณาอนุพันธ์อันดับสอง
\[ \dx[2] f(x) = \dx f\,'(x) = \limh \frac{f\,'(x+h) - f\,'(x)}{h}\]
ซึ่ง $\dx f\,'(x)$ คือการหาความเปลี่ยนแปลงของฟังก์ชัน $f\,'(x)$ เทียบกับค่า $x$ โดย นั่นคือ 
\begin{block}{}
อนุพันธ์อันดับสอง คือการอัตราความเปลี่ยนแปลงของความเปลี่ยนแปลงของฟังก์ชัน
\end{block}
\end{frame}

\begin{frame}
ถ้า $f\,''(x) > 0$ แสดงว่าฟังก์ชัน $f\,'(x)$ เป็นฟังก์ชันเพิ่ม นั่นคืออนุพันธ์ของกราฟ $ y = f(x)$ เพิ่มขึ้น โดยดูได้จากกราฟดังต่อไปนี้ ซึ่งมีชื่อเรียกว่าเว้าบน (concave up)
\addfig{0.6}{diff_applied_concave_up}{ตัวอย่างกราฟ $y=f(x)$ ที่อนุพันธ์เพิ่มขึ้น}{}
\end{frame}

\begin{frame}
ถ้า $f\,''(x) < 0$ แสดงว่าฟังก์ชัน $f\,'(x)$ เป็นฟังก์ชันลด นั่นคืออนุพันธ์ของกราฟ $y = f(x)$ ลดลง โดยดูได้จากกราฟดังต่อไปนี้ ซึ่งมีชื่อเรียกว่าเว้าล่าง (concave down)
\addfig{0.6}{diff_applied_concave_down}{ตัวอย่างกราฟ $y=f(x)$ ที่อนุพันธ์ลดลง}{}
\end{frame}

\begin{frame}[t]
ถ้า $f\,''(x) = 0$ เรายังไม่สามารถสรุปได้ทันทีว่าพฤติกรรมของเส้นโค้ง ณ จุดนั้นเป็นอย่างไร ต้องพิจารณาจากกราฟของ $f\,'(x)$ โดยดูค่าใกล้ ๆ นั่นคือ $f\,'(x-c)$ และ $f\,'(x+c)$ เมื่อ $c$ เป็นค่าคงตัวที่เล็ก
\begin{slide}
\begin{itemize}
\item ถ้า $f\,'(x-c)$ และ $f\,'(x+c)$ มีเครื่องหมายเดียวกัน (เป็นบวกทั้งคู่ หรือ เป็นลบทั้งคู่) แสดงว่าโค้งทั้งสองด้านเป็นโค้งเว้าแบบเดียวกัน นั่นคือเว้นบนทั้งคู่หรือเว้าล่างทั้งคู่ ฉะนั้นเส้นโค้ง ณ จุดที่ $f\,''(x) = 0$ นั้นเป็นเส้นโค้งเว้าแบบเดียวกันด้วย
\item ถ้า $f\,'(x-c)$ และ $f\,'(x+c)$ มีเครืองหมายต่างกัน (เป็นลบตัวหนึ่ง เป็นบวกตัวหนึ่ง) แสดงว่าโค้งทั้สองด้านเป็นโค้งเว้าคนละแบบกัน เราจะเรียกจุด $(x,f(x))$ นั้นว่า \textbf{จุดเปลี่ยนเว้า}
\end{itemize}
\end{slide}
\end{frame}

\begin{frame}
\begin{columns}
\begin{column}{0.5\textwidth}
\addfig{1}{diff_applied_inflection1}{จุดเปลี่ยนเว้า จากเว้าลงเป็นเว้าขึ้น}{}
\end{column}
\begin{column}{0.5\textwidth}
\addfig{1}{diff_applied_inflection2}{จุดเปลี่ยนเว้า จากเว้าขึ้นเป็นเว้าลง}{}
\end{column}
\end{columns}
\end{frame}

\begin{frame}
\frametitle{นิยามจุดเปลี่ยนโค้ง}
\begin{defn}{}{}
\begin{itemize}
\item $f(x)$ เป็นฟังก์ชัน\textbf{เว้าบน} (concave upward) ในช่วง $[a,b]$ ถ้า $f\,''(x) > 0$ สำหรับทุกค่า $x$ ในช่วง $[a,b]$
\item $f(x)$ เป็นฟังก์ชัน\textbf{เว้าล่าง} (concave downward) ในช่วง $[a,b]$ ถ้า $f\,''(x) < 0$ สำหรับทุกค่า $x$ ในช่วง $[a,b]$
\item \textbf{จุดเปลี่ยนความเว้า} (point of inflection) คือจุดที่เส้นโค้งเปลี่ยนจากเว้าบนเป็นเว้าล่าง หรือเว้าล่างเป็นเว้าบน ซึ่ง $f\,''(x) = 0$
\end{itemize}
\end{defn}
\end{frame}


\begin{frame}[t]
\begin{ex}
จงหาว่าฟังก์ชัน $f(x) = x^3-6x^2$ เป็นฟังก์ชันเว้าบนในช่วงใด และเป็นฟังก์ชันเว้าล่างในช่วงใด พร้อมหาจุดเปลี่ยนเว้าด้วย (ถ้ามี)
\end{ex}
\begin{sol}
พิจารณาอนุพันธ์อันดับสอง 
\[ f\,'(x) = \dx f(x) = 3x^2 - 12x \qquad \then \qquad f\,''(x) = \dx f\,'(x) = 6x - 12 \]
\begin{itemize}
\item $f\,''(x) = 6x-12 > 0$ เมื่อ $x > 2$ ฉะนั้น $f(x)$ เป็นฟังก์ชันเว้าบนในช่วง $(2, \infty)$
\item $f\,''(x) = 6x-12 < 0$ เมื่อ $x < 2$ ฉะนั้น $f(x)$ เป็นฟังก์ชันเว้าล่างในช่วง $(-\infty, 2)$
\item ค่า $x=2$ เป็นค่าที่ $f(x)$ เปลี่ยนจากฟังก์ชันเว้าบน เป็นฟังก์ชันเว้าล่าง ดังนั้น $f(x)$ มีจุดเปลี่ยนเว้าเมื่อ $x=2$ ซึ่งเป็นจุด $(2, f(x)) = (2, -16)$
\end{itemize}
\end{sol}
\end{frame}

%%%%%%%%%%%%%%
\section[ค่าสุดขีด]{ค่าสุดขีดของฟังก์ชัน}
\frame{\sectionpage}
%%%%%%%%%%%%%%

\begin{frame}
\frametitle{การประยุกต์การหาค่าสุดขีดสัมพัทธ์ของฟังก์ชัน}
ในการแก้ปัญหาทางวิศวกรรมหรือธุรกิจ เรามักจะเจอโจทย์ปัญหาต่อไปนี้
\begin{itemize}
\item หาตำแหน่งสูงสุดของลูกบอลที่เราโยนขึ้นฟ้า
\item หาระยะทางที่ไกลที่สุด ที่เกิดจากการดีดสปริง
\item หาจำนวนสินค้าที่ต้องผลิต ถ้าต้องการให้ได้กำไรมากที่สุด
\end{itemize}
ซึ่งแปลงเป็นโจทย์ทางคณิตศาสตร์ได้เป็น
\begin{block}{}
กำหนดฟังก์ชัน $f(x)$ ในช่วง $[a,b]$ จงหาค่า $x$ ที่ทำให้ $f(x)$ มากที่สุด หรือ น้อยที่สุด
\end{block}
ถ้าฟังก์ชัน $f(x)$ ไม่ซับซ้อนนัก เราสามารถใช้อนุพันธ์ในการหาค่าที่ดีที่สุดได้
\end{frame}

\begin{frame}
\addfig{0.8}{diff_applied_extrema_intro}{การจัดการสินค้า จำเป็นต้องคำนวณต้นทุนให้คุ้มค่าที่สุด ภายใต้ตัวแปรต่าง ๆ \href{https://www.pexels.com/photo/birds-eye-view-photo-of-freight-containers-2226458/}{(ที่มา)}}{}
\end{frame}


\begin{frame}
\frametitle{ค่าสุดขีดสัมบูรณ์และสัมพัทธ์}
\begin{defn}{}{}
กำหนดให้ $f$ เป็นฟังก์ชันนิยามบนช่วงปิด $[a,b]$ และ $c \in [a,b]$
\begin{enumerate}
\item $f$ มี\textbf{ค่าสูงสุดสัมบูรณ์} ที่ $x=c$ \pa ก็ต่อเมื่อ $f(c)$ เป็นค่ามากที่สุดในช่วง $[a,b]$ \pa (นั่นคือ $f(c) \geq f(x)$ สำหรับทุกค่า $x \in [a,b]$) 
\item $f$ มี\textbf{ค่าต่ำสุดสัมบูรณ์} ที่ $x=c$ \pa ก็ต่อเมื่อ $f(c)$ เป็นค่าน้อยที่สุดในช่วง $[a,b]$ \pa (นั่นคือ $f(c) \leq f(x)$ สำหรับทุกค่า $x \in [a,b]$)
\item $f$ มีค่าสูงสุด\textbf{สัมพัทธ์} ที่ $x=c$ \pa ก็ต่อเมื่อ $f(c)$ เป็นค่ามากที่สุดในช่วง $[a_0,b_0]$ สำหรับค่า $a_0, b_0$ บางค่าที่ $a_0 < c < b_0$  
\item $f$ มีค่าต่ำสุด\textbf{สัมพัทธ์} ที่ $x=c$ \pa ก็ต่อเมื่อ $f(c)$ เป็นค่าน้อยที่สุดในช่วง $[a_0,b_0]$ สำหรับค่า $a_0, b_0$ บางค่าที่ $a_0 < c < b_0$  
\end{enumerate}
\end{defn}
\end{frame}


\begin{frame}[t]
\begin{ex}
\begin{columns}
\begin{column}{0.7\textwidth}
จงหาว่าจุดต่าง ๆ ในกราฟต่อไปนี้ เป็นจุดประเภทใด ถ้ากำหนดให้ $f$ เป็นฟังก์ชันที่นิยามในช่วงของกราฟ
\end{column}
\begin{column}{0.3\textwidth}
\addfig{1}{diff_applied_extrema_ex1}{กราฟ $y = f(x)$}{}
\end{column}
\end{columns}
\end{ex}
\begin{sol}
\begin{itemize}
\item จุด $C$ เป็นจุดต่ำสุดสัมบูรณ์
\item จุด $D$ เป็นจุดสูงสุดสัมบูรณ์
\item จุด $B$ และ $D$ เป็นจุดสูงสุดสัมพัทธ์
\item จุด $A$ และ $C$ เป็นจุดต่ำสุดสัมพัทธ์
\end{itemize}
\end{sol}
\end{frame}

\begin{frame}
\frametitle{หมายเหตุ}
\begin{remark}{}
\begin{itemize}
\item $f$ มี\textbf{ค่า}สูงสุดสัมบูรณ์มีได้ค่าเดียว หรืออาจจะไม่มีเลยก็ได้
\item $f$ มี\textbf{ค่า}ต่ำสุดสัมบูรณ์มีได้ค่าเดียว หรืออาจจะไม่มีเลยก็ได้
\item $f$ มี\textbf{จุด}สูงสุดสัมบูรณ์และจุดต่ำสุดสัมบูรณ์ได้หลายจุด เช่น ฟังก์ชัน $f(x) = \cos x$ มีค่าสูงสุดสัมพัทธ์เท่ากับ $1$ แต่มีจุดสูงสุดสัมพัทธ์เมื่อ $x = n \pi$ สำหรับทุกจำนวนเต็ม $n$
\end{itemize}
\end{remark}
\end{frame}

\begin{frame}
\addfig{1}{diff_applied_extrema_cos}{ฟังก์ชัน $f(x) = \cos x$ มีค่าสูงสุดสัมพัทธ์และค่าต่ำสุดสัมพัทธ์ที่ค่า $x$ หลายค่า}{}
\end{frame}

\begin{frame}
\frametitle{จุดวิกฤต}
คำถามต่อไปคือ เราจะทราบได้อย่างไรว่าฟังก์ชัน $f$ มีค่าสูงสุดสัมพัทธ์และต่ำสุดสัมพัทธ์ที่ใด
\begin{defn}{}{}
\textbf{ค่าวิกฤต} ของฟังก์ชัน $f(x)$ \pa คือค่า $c$ ที่ทำให้
\begin{itemize}
\item $f\,'(c) = 0$ หรือ \pa 
\item $f\,'(c)$ หาค่าไม่ได้
\end{itemize}
ถ้า $c$ เป็นค่าวิกฤตของฟังก์ชัน แล้วเราเรียกจุด $(c,f(c))$ บนกราฟว่า\textbf{จุดวิกฤต}
\end{defn}
\end{frame}


\begin{frame}[t]
\begin{ex}
จงหาค่าวิกฤตทั้งหมดของฟังก์ชันต่อไปนี้
\begin{tasks}(2)
\task $f(x) = 3x$
\task $f(x) = x^2-4x$
\end{tasks}
\end{ex}
\begin{sol}
\begin{tasks}(1)
\task ฟังก์ชัน $f(x) = 3x$ \pa หาอนุพันธ์ได้ 
\[f\,'(x) = 3\]
\pa แสดงว่าเราหาค่า $c$ ที่ทำให้ $f\,'(c)=0$ ไม่ได้ \pa นั่นคือ $f(x)$ ไม่มีค่าวิกฤต \pa
\task ฟังก์ชัน $f(x) = x^2-4x$ \pa หาอนุพันธ์ได้ 
\[f\,'(x) = 2x-4\] \pa 
ซึ่งมีค่าเท่ากับ $0$ เมื่อ $x=2$ \pa ฉะนั้น $f(x)$ มีค่าวิกฤต $x=2$ \pa
\end{tasks}
\end{sol}
\end{frame}



\begin{frame}[t]
\begin{ex}
จงหาค่าวิกฤตทั้งหมดของฟังก์ชันต่อไปนี้
\begin{tasks}(2)
\task $f(x) = x^3-3x$
\task $f(x) = |x|$
\end{tasks}
\end{ex}
\begin{sol}
\begin{tasks}(1)
\task ฟังก์ชัน $f(x) = x^3-3x$ \pa หาอนุพันธ์ได้ 
\[f\,'(x) = 3x^2-3\] \pa 
ซึ่งมีค่าเท่ากับ $0$ เมื่อ $x =-1, 1$ \pa ฉะนั้นเ $f(x)$ มีค่าวิกฤตสองค่า คือเมื่อ $x = -1$ และ $x=1$ \pa
\task ฟังก์ชัน $f(x) = |x|$ \pa หาอนุพันธ์ได้ $f\,'(x) = -1$ เมื่อ $x<0$, $f\,'(x) = 1$ เมื่อ $x> 0$ และหาค่าไม่ได้เมื่อ $x=0$ ฉะนั้น $f(x)$ มีค่าวิกฤต $x=0$ 
\end{tasks}
\end{sol}
\end{frame}

\begin{frame}
\frametitle{การทดสอบจุดวิกฤต}
เมื่อได้ค่าวิกฤตและจุดวิกฤตแล้ว เราต้องพิจารณาต่อว่าจุดนั้นเป็นจุดประเภทใด สังเกตว่าถ้า $f\,'(c) = 0$ แล้วเรายังไม่สามารถสรุปได้ทันทีว่าเป็นจุดประเภทใด แต่เราสามารถใช้ความเว้าของกราฟมาวิเคราะห์ได้ทันทีดังนี้
\begin{itemize}
\item จุดสูงสุดสัมพัทธ์ จะมีรูปเว้าบน นั่นคือค่าของ $f\,'(x)$ จะเปลี่ยนจากค่าบวก เป็นศูนย์ และเป็นค่าลบ นั่นคือ $f\,''(x)$ เป็นลบ
\item จุดต่ำสุดสัมพัทธ์ จะมีรูปเว้าล่าง นั่นคือค่าของ $f\,'(x)$ จะเปลี่ยนจากค่าลบ เป็นศูนย์ และเป็นค่าบวก นั่นคือ $f\,''(x)$ เป็นบวก
\end{itemize}
\end{frame}

\begin{frame}
\frametitle{การทดสอบจุดวิกฤต}
\begin{thm}{}{}
ให้ $f(x)$ เป็นฟังก์ชันที่หาอนุพันธ์และอนุพันธ์อันดับสองได้บนช่วง $[a,b]$ \pa และ $c$ เป็นค่าวิกฤตแล้ว\pa
\begin{enumerate}
\item ถ้า $f\,''(c) < 0$ แล้ว $f(c)$ เป็นค่าสูงสุดสัมพัทธ์\pa
\item ถ้า $f\,''(c) > 0$ แล้ว $f(c)$ เป็นค่าต่ำสุดสัมพัทธ์\pa
\item ถ้า $f\,''(c) = 0$ แล้วเรายังไม่สามารถให้ข้อสรุปได้ ต้องพิจารณาจาก $f\,'(x)$ อีกที
\end{enumerate}
\end{thm}
\end{frame}


\begin{frame}[t]
\begin{ex}
จงหาจุดวิกฤตทั้งหมดของฟังก์ชัน $f(x) = x^2-4x$ พร้อมทดสอบว่าเป็นจุดประเภทใด 
\end{ex}\pa 
\begin{sol}
อนุพันธ์ของฟังก์ชันมีค่าเท่ากับ 
\[f\,'(x) = 2x-4\] \pa 
ซึ่งมีจุดวิกฤตเมื่อ $2x-4 = 0$ นั่นคือ $x_0 = 2$ \pa อนุพันธ์อันดับสองคือ 
\[f\,''(x) = 2\] \pa 
ทำให้ $f\,''(x_0) = 2 > 0$ \pa ดังนั้นจุดวิกฤต $x_0 = 2$ คือจุดต่ำสุดสัมพัทธ์
\end{sol}
\end{frame}


\begin{frame}[t]
\begin{ex}
จงหาจุดวิกฤตทั้งหมดของฟังก์ชัน $f(x) = -2x^2-8x+10$ พร้อมทดสอบว่าเป็นจุดประเภทใด 
\end{ex}\pa 
\begin{sol}
อนุพันธ์ของฟังก์ชันมีค่าเท่ากับ 
\[f\,'(x) = -4x-8\] \pa 
ซึ่งมีจุดวิกฤตเมื่อ $-4x-8 = 0$ นั่นคือ $x_0 = -2$ \pa อนุพันธ์อันดับสองคือ 
\[f\,''(x) = -4\] \pa 
ทำให้ $f\,''(x_0) = -4 < 0$ \pa ดังนั้นจุดวิกฤต $x_0 = -2$ คือจุดสูงสุดสัมพัทธ์
\end{sol}
\end{frame}


\begin{frame}[t]
\begin{ex}
จงหาจุดวิกฤตทั้งหมดของฟังก์ชัน $f(x) = x^3-3x$ พร้อมทดสอบว่าเป็นจุดประเภทใด 
\end{ex} 
\begin{sol}
อนุพันธ์ของฟังก์ชันมีค่าเท่ากับ 
\[f\,'(x) = 3x^2-3\] \pa 
ซึ่งมีจุดวิกฤตเมื่อ $3x^2-3 = 0$ นั่นคือ $x_0 = -1, 1$ \pa รวม 2 จุด โดยเราต้องวิเคราะห์แต่ละจุดแยกกัน

อนุพันธ์อันดับสองคือ 
\[f\,''(x) = 6x\] \pa
\begin{itemize}
\item สำหรับจุดวิกฤต $x_0 = -1$ \pa เราทราบว่า $f\,''(-1) = -6 < 0$\pa  ดังนั้นจุดวิกฤต $x_0 = -1$ คือจุดสูงสุดสัมพัทธ์\pa 
\item สำหรับจุดวิกฤต $x_0 = 1$\pa  เราทราบว่า $f\,''(-1) = 6 > 0$ \pa ดังนั้นจุดวิกฤต $x_0 = 1$ คือจุดต่ำสุดสัมพัทธ์
\end{itemize}
\end{sol}
\end{frame}


\begin{frame}[t]
\begin{ex}
จงหาจุดวิกฤตทั้งหมดของฟังก์ชัน $f(x) = x^6$ พร้อมทดสอบว่าเป็นจุดประเภทใด 
\end{ex} 
\begin{sol}
อนุพันธ์ของฟังก์ชันมีค่าเท่ากับ 
\[f\,'(x) = 6x^5\] \pa 
ซึ่งมีจุดวิกฤตเมื่อ $x_0 = 0$ \pa อนุพันธ์อันดับสองคือ 
\[f\,''(x) = 30x^4\] \pa 
ทำให้ $f\,''(x_0) = 0$ \pa ดังนั้นเราไม่สามารถสรุปได้ว่าจุดวิกฤต $c = 0$ เป็นจุดประเภทใด
\end{sol}
\end{frame}

%%%%%%%%%%%%%%
\section[โจทย์ปัญหาค่าสุดขีด]{โจทย์ปัญหาเรื่องค่าสุดขีดของฟังก์ชัน}
\frame{\sectionpage}
%%%%%%%%%%%%%%

\begin{frame}
\frametitle{โจทย์ปัญหาค่าสุดขีด}
ถ้าโจทย์กำหนดฟังก์ชันมาให้แล้ว เราสามารถหาจุดวิกฤตและทดสอบได้ทันที แต่ปกติแล้วโจทย์ปัญหาที่เราจะพบเจอจะอยู่ในรูปแบบของสถานการณ์ ที่เราต้องสร้างสมการทางคณิตศาสตร์ก่อนถึงจะแก้ได้ โดยมีรูปแบบการแก้โจทย์ดังนี้
\begin{enumerate}
\item กำหนดตัวแปรต่าง ๆ
\item สร้างสมการความสัมพันธ์ โดยจัดรูปให้เหลือเพียงตัวแปรเดียว
\item ทดสอบค่าสุดขีด
\end{enumerate}
ดังจะได้เห็นตามตัวอย่างต่อไปนี้

\begin{remark}{}
ถ้าโจทย์กำหนดช่วง $[a,b]$ ของฟังก์ชันมาให้ ซึ่งเป็นช่วงปิด เราต้องพิจารณาค่าขอบ $x=a$ และ $x=b$ เนื่องจาก ณ จุดนั้น $f\,'(x)$ หาค่าไม่ได้ จึงนับว่าเป็นจุดวิกฤตด้วย
\end{remark}
\end{frame}


\begin{frame}[t]
\begin{ex}
มีเชือกยาว 20 เซนติเมตร จะสร้างเป็นรูปสี่เหลี่ยมมุมฉากให้มีพื้นที่มากที่สุดได้เท่าไหร่ และด้านแต่ละด้านจะมีความยาวเท่าไหร่
\end{ex} 
\begin{columns}
\begin{column}{0.7\textwidth}
\begin{sol} กำหนดตัวแปรดังนี้
\begin{itemize}
\item $x$ เป็นความกว้าง 
\item $y$ เป็นความยาว 
\item $A$ เป็นพื้นที่ 
\end{itemize}
จากข้อจำกัดของความยาวเชือกเราทราบว่า
\[x+y+x+y = 20\] 
เมื่อจัดรูปจะได้
\[x+y = 10\]
\end{sol}
\end{column}
\begin{column}{0.3\textwidth}
\addfig{0.8}{diff_applied_extrema_problem1}{สี่เหลี่ยมมุมฉากความกว้าง $x$ ความยาว $y$}{}
\end{column}
\end{columns}
\end{frame}

\begin{frame}
\begin{sol} (ต่อ)
แต่เนื่องจากเราต้องใช้ตัวแปรเดียว ฉะนั้นเราจะเขียน $y = 10-x$ จากนั้น กำหนดฟังก์ชันพื้นที่สี่เหลี่ยมมุมฉากเป็น
\[A = xy = x(10-x) = 10x - x^2\]
นั่นคือ $A(x) = 10x - x^2$ เราหาจุดวิกฤตได้ด้วยการหาอนุพันธ์ 
\[A'(x) = 10 - 2x\]
ซึ่งมีค่าเท่ากับ 0 เมื่อ $x=5$ และเนื่องจาก $A''(x) = -2$ ซึ่งน้อยกว่าศูนย์ ฉะนั้นจุดนี้เป็นจุดสูงสุดสัมพัทธ์ เราจึงสรุปได้ว่า สี่เหลี่ยมจะมีพื้นที่มากที่สุดเมื่อมีความกว้าง 5 เซนติเมตร และยาว $10-w=10-5=5$ เซนติเมตร และมีพื้นที่เท่ากับ
\[A(5) = 5(10-5) = 25\]
โดยมีหน่วยเป็นตารางเซนติเมตร
\end{sol}
\end{frame}


\begin{frame}[t]
\begin{ex}
บริษัทแห่งนึงต้องการสร้างกระป๋องน้ำอัดลมเป็นทรงกระบอกที่จุน้ำได้ 250 ลูกบาศก์มิลลิเมตร โดยต้องการให้พื้นที่ผิวของวัสดุมีค่ามากที่สุด จะต้องผลิตกระป๋องที่มีความสูงและรัศมีเท่าใด
\end{ex} 
\begin{columns}
\begin{column}{0.7\textwidth}
\begin{sol} กำหนดตัวแปรดังนี้
\begin{itemize}
\item $x$ เป็นรัศมีของฝากระป๋อง
\item $y$ เป็นความสูง
\item $S$ เป็นพื้นที่ผิวทั้งหมด
\end{itemize}
เราคำนวณปริมาตรของกระป๋องน้ำได้เป็น
\[y(\pi x^2) = 250\] 
เราเขียน $y$ ในรูปของ $x$ ได้เป็น $y = \frac{250}{\pi x^2}$
\end{sol}
\end{column}
\begin{column}{0.3\textwidth}
\addfig{0.8}{diff_applied_extrema_problem2-1}{ทรงกระบอกความสูง $y$ และรัศมี $x$}{}
\end{column}
\end{columns}
\end{frame}

\begin{frame}
\begin{sol} (ต่อ)

พื้นที่ผิวประกอบด้วยวงกลมที่ฝาและก้นของกระป๋อง และสี่เหลี่ยมผืนผ้าด้านข้างกระป๋อง ดังภาพต่อไปนี้
\addfig{0.6}{diff_applied_extrema_problem2-2}{พื้นที่ผิวของทรงกระบอก}{}
ซึ่งรวมกันได้เป็น
\[S = 2(\pi x^2) + (2\pi x)(y) = 2\pi x^2 + 2\pi x \left( \frac{250}{\pi x^2}\right) = 2\pi x^2 + \frac{500}{x}\]
นั่นคือ กำหนดให้ $S(x) = 2\pi x^2 + \tfrac{500}{x}$ 
\end{sol}
\end{frame}

\begin{frame}
\begin{sol} (ต่อ)
เราหาจุดวิกฤตได้ด้วยการหาอนุพันธ์ 
\begin{eq*}
S'(x) = 4\pi x - \frac{500}{x^2} &=& 0 \\
4\pi x &=& \frac{500}{x^2} \\
x^3 &=& \frac{125}{\pi}\\
x &=& \frac{5}{\sqrt[3]{\pi}}
\end{eq*}

จากนั้นเราสามารถทดสอบจุดวิกฤตนี้ได้ด้วยการหาอนุพันธ์อันดับสอง
\[S''(x) = 4\pi + \frac{1500}{x^3}\]
โดยเมื่อแทนค่า $x = \frac{5}{\sqrt[4]{\pi}}$ ลงไป จะเห็นได้ว่าคำตอบจะเป็นค่าบวกเนื่องจากทุกพจน์เป็นพจน์บวก ฉะนั้นจุดนี้เป็นจุดต่ำสุดสัมพัทธ์ 
\end{sol}
\end{frame}

\begin{frame}
\begin{sol} (ต่อ)
เราจึงสรุปได้ว่า กระป๋องนี้จะมีพื้นที่ผิวมากที่สุดเมื่อ
\begin{itemize}
\item รัศมี $x=\frac{5}{\sqrt[3]{\pi}} \approx 3.414$ มิลลิเมตร 
\item ความสูง $y = \frac{250}{\pi x^2} = \frac{25}{\sqrt[4]{\pi}} \approx 6.828$ มิลลิเมตร
\item พื้นที่ผิวเท่ากับ $S(x) = S\left(\frac{5}{\sqrt[3]{\pi}}\right) \approx 219.689$ ตารางมิลลิเมตร
\end{itemize}
\end{sol}
\end{frame}